% SGA 6, Exposición IX: fragmento español alineado con la fuente.
% Alcance: apertura y sección 1, líneas 11991--12198 inclusive.
% Fuente francesa: sga6_fr_r1/src/fr.tex.
% SHA-256 de la fuente francesa completa:
% 0D2158B9FE3B466309A92165727A0B191422E05FBAA0DACD5806E8322C54BA1A.
\providecommand{\sisi}[2]{#2}

\section*{Exposición IX. Algunos cálculos de grupos $K_\bullet$}
\begin{center}
por P. Berthelot
\end{center}

Esta exposición no corresponde a ninguna exposición oral y no será utilizada en la continuación del seminario. Contiene cierto número de cálculos de grupos $K_\bullet$ para diversas fibraciones: fibrados vectoriales, fibrados proyectivos, torsores bajo un grupo lineal, etc.\ldots{} En toda la exposición se supondrá noetheriano el esquema base, y la notación $K_\bullet(X)$ designará el grupo de Grothendieck de los haces coherentes sobre un esquema $X$. Recordemos que, si $X$ es noetheriano, el homomorfismo del $K$ de la categoría de los $\mathcal O_X$-módulos coherentes al $K$ de la categoría de los complejos pseudocoherentes sobre $X$ es un isomorfismo (IV 2.4).

Los resultados expuestos aquí son análogos a los que se obtienen para el anillo de Chow y que se encontrarán en las exposiciones de Grothendieck en el Seminario Chevalley de 1958. Se plantea asimismo la cuestión de determinar en qué medida pueden obtenerse enunciados del mismo tipo para el graduado asociado a $K_\bullet$ respecto de la filtración definida por la dimensión del soporte (cf. X). No parece que se conozca la respuesta, ni siquiera para la proposición 1.1 (sucesión exacta de homotopía), que es una de las propiedades más utilizadas en los cálculos siguientes.

\subsection*{1. Fibrados vectoriales}

\textbf{Proposición 1.1 (sucesión exacta de homotopía).} Sean $X$ un esquema noetheriano, $X'$ un subesquema cerrado de $X$, $U$ el abierto complementario de $X'$ en $X$, $i$ la inmersión de $X'$ en $X$, y $j$ la inmersión de $U$ en $X$. Entonces la sucesión
\[
K_\bullet(X')\xrightarrow{i_*}K_\bullet(X)\xrightarrow{j^*}K_\bullet(U)\longrightarrow 0
\]
es exacta.

Si $F'$ es un $\mathcal O_{X'}$-módulo coherente, $i_*(F')$ es un $\mathcal O_X$-módulo cuyo soporte está contenido en $X'$; por consiguiente, su restricción a $U$ es nula, y $j^*\circ i_*(\cl(F'))=0$. Como los elementos de la forma $\cl(F')$ engendran $K_\bullet(X')$, se sigue que $j^*\circ i_*=0$. Se deduce un homomorfismo
\[
 \rho:K_\bullet(X)/\operatorname{Im}K_\bullet(X')\longrightarrow K_\bullet(U).
\]
Vamos a definir un homomorfismo inverso de $\rho$. Sea $G$ un $\mathcal O_U$-módulo coherente. Como $X$ es noetheriano, existe un $\mathcal O_X$-módulo coherente $F$ cuya restricción a $U$ es isomorfa a $G$. Si $F'$ es un segundo $\mathcal O_X$-módulo coherente cuya restricción a $U$ es isomorfa a $G$, entonces se tiene en $K_\bullet(X)$:
\[
 \cl(F)-\cl(F')\in \operatorname{Im}K_\bullet(X').
\]
Para verlo, consideramos $G$ como un subhaz de $j^*(F\oplus F')$ mediante la aplicación diagonal; $G$ se prolonga entonces a un subhaz $F''$ de $F\oplus F'$, y quedamos reducidos a demostrar que
\[
 \cl(F)-\cl(F'')\in \operatorname{Im}K_\bullet(X'),
\]
y el análogo para $F'$. Sean entonces $H$ y $K$ el núcleo y el conúcleo del homomorfismo compuesto
\[
F''\longrightarrow F\oplus F'\longrightarrow F ;
\]
este homomorfismo se reduce a la identidad de $G$ sobre $U$, y por tanto $H$ y $K$ tienen soporte contenido en $X'$. Si $I$ es el ideal de $X'$ en $X$, se sigue que existen enteros $m$ y $n$ tales que $I^m\cdot H=0$, $I^n\cdot K=0$ (EGA I 9.3.4). Por tanto, se pueden descomponer $H$ y $K$ en extensiones de módulos anulados por $I$, y por ello de $\mathcal O_{X'}$-módulos, y se tiene
\[
\cl(H)\in \operatorname{Im}K_\bullet(X'),\qquad
\cl(K)\in \operatorname{Im}K_\bullet(X'),
\]
de donde resulta la afirmación. Definimos así una función sobre la categoría de los $\mathcal O_U$-módulos coherentes con valores en $K_\bullet(X)/\operatorname{Im}K_\bullet(X')$, asociando a todo $\mathcal O_U$-módulo coherente la clase de una de sus prolongaciones sobre $X$.

Esta función es aditiva. En efecto, sea
\[
0\longrightarrow G'\longrightarrow G\longrightarrow G''\longrightarrow 0
\]
una sucesión exacta de $\mathcal O_U$-módulos coherentes. Sea $F$ una prolongación de $G$ sobre $X$; existe entonces un submódulo $F'$ de $F$ que prolonga $G'$, y $F/F'$ prolonga $G''$. La función se factoriza, pues, por $K_\bullet(U)$, lo que define el homomorfismo buscado $K_\bullet(U)\to K_\bullet(X)/\operatorname{Im}K_\bullet(X')$. Se comprueba inmediatamente que es inverso de $\rho$.

\textbf{Proposición 1.2.} Sean $X,S$ dos esquemas noetherianos, y $f$ un morfismo plano de $X$ a $S$. Se da una familia de elementos $(\xi_i)_{i\in I}$ de $K^\circ(X)$, y se considera el homomorfismo
\[
\varphi:(K_\bullet(S))^{(I)}\longrightarrow K_\bullet(X)
\]
definido por
\[
\varphi((x_i)_{i\in I})=\sum_{i\in I}\xi_i\cdot f^*(x_i),
\]
donde $f^*$ es el homomorfismo de «~imagen inversa~» de $K_\bullet(S)$ en $K_\bullet(X)$, definido según IV 2.12, y donde $K_\bullet(X)$ se considera como $K^\circ(X)$-módulo por IV 2.10. Se supone que, para todo punto $s\in S$, las imágenes de los $\xi_i$ en $K_\bullet(X_s)$ por el homomorfismo
\[
K^\circ(X)\longrightarrow K^\circ(X_s)\longrightarrow K_\bullet(X_s)
\]
engendran $K_\bullet(X_s)$ como grupo abeliano. Entonces $\varphi$ es sobreyectivo.

Sean $S'$ un subesquema cerrado de $S$, y $U$ el abierto complementario de $S'$ en $S$. Denotamos por $X_U$ la restricción de $X$ sobre $U$, y por $X_{S'}$ la restricción de $X$ sobre $S'$. Finalmente, denotamos todavía por $\xi_i$ las imágenes de los $\xi_i$ en $K^\circ(X_{S'})$ y $K^\circ(X_U)$, y escribimos
\[
\varphi_{S'}:K_\bullet(S')^{(I)}\to K_\bullet(X_{S'}),\qquad
\varphi_U:K_\bullet(U)^{(I)}\to K_\bullet(X_U)
\]
para los homomorfismos definidos del mismo modo que $\varphi$.

Consideremos el diagrama
\[
\begin{tikzcd}[column sep=tiny]
K_\bullet(S')^{(I)} \arrow[r,"i_*"] \arrow[d,"\varphi_{S'}"'] & K_\bullet(S)^{(I)} \arrow[r,"j^*"] \arrow[d,"\varphi"] & K_\bullet(U)^{(I)} \arrow[r] \arrow[d,"\varphi_U"] & 0\\
K_\bullet(X_{S'}) \arrow[r,"u_*"] & K_\bullet(X) \arrow[r,"v^*"] & K_\bullet(X_U) \arrow[r] & 0,
\end{tikzcd}\tag{1.2.1}
\]
donde $i$ y $j$ designan las inmersiones de $S'$ y $U$ en $S$, y $u$ y $v$ las inmersiones de $X_{S'}$ y $X_U$ en $X$. Demostremos que, si $\varphi_{S'}$ y $\varphi_U$ son sobreyectivos, lo mismo vale para $\varphi$.

Las dos líneas horizontales del diagrama son exactas por 1.1. Se ve entonces que, para demostrar que $\varphi$ es sobreyectivo, basta demostrar que los dos cuadrados son conmutativos. Sea $(x'_i)_{i\in I}\in K_\bullet(S')^{(I)}$~; se tiene
\[
\varphi\circ i_*((x'_i))=\varphi((i_*(x'_i)))=\sum_i \xi_i\cdot f^*(i_*(x'_i));
\]
\[
u_*\circ\varphi_{S'}((x'_i))=u_*\left(\sum_i \xi_i\cdot f^*_{S'}(x'_i)\right)=\sum_i \xi_i\cdot u_*(f^*_{S'}(x'_i)),
\]
donde la última igualdad resulta de la fórmula de proyección IV 2.11.1.2. La igualdad $\varphi\circ i_*=u_*\circ\varphi_{S'}$ resulta, pues, de la igualdad $f^*\circ i_*=u_*\circ f^*_{S'}$, válida por IV 3.1.1, pues $i$ es propio y $f$ plano. Sea $(x_i)_{i\in I}\in K_\bullet(S)^{(I)}$~; se tiene
\[
\varphi_U\circ j^*((x_i))=\sum_i \xi_i\cdot f_U^*\circ j^*(x_i);
\]
\[
v^*\circ\varphi((x_i))=v^*\left(\sum_i \xi_i\cdot f^*(x_i)\right)=\sum_i \xi_i\cdot v^*\circ f^*(x_i)=\sum_i \xi_i\cdot f_U^*\circ j^*(x_i),
\]
donde la igualdad central resulta de IV 2.12.2.

Este resultado permite demostrar 1.2 por inducción noetheriana sobre el conjunto de los subesquemas cerrados de $S$. Si suponemos que, para todo subesquema cerrado $S'$ de $S$ distinto de $S$, $\varphi_{S'}$ es sobreyectivo, el razonamiento anterior muestra que podemos reducirnos a demostrar 1.2 para un esquema base irreducible y afín, reemplazando si es necesario $S$ por un abierto afín irreducible. Además, el diagrama
\[
\begin{tikzcd}
K_\bullet(S_{\red})^{(I)} \arrow[r] \arrow[d,"\varphi_{S_{\red}}"'] & K_\bullet(S)^{(I)} \arrow[d,"\varphi"]\\
K_\bullet(X_{S_{\red}}) \arrow[r] & K_\bullet(X)
\end{tikzcd}
\]
donde $K_\bullet(X_{S_{\red}})\to K_\bullet(X)$ es sobreyectivo por 1.1, muestra que también podemos suponer $S$ reducido.

Nos situamos, pues, sobre un esquema base íntegro y afín. Ponemos $S=\Spec(A)$ y denotamos por $K$ el cuerpo de fracciones de $A$; si $s$ es el punto genérico de $S$, entonces $\mathcal O_{S,s}=K$. Sea $y\in K_\bullet(X)$. Por hipótesis, $K_\bullet(X_s)$ está engendrado por las imágenes de los $\xi_i$; si $y_s$ es la imagen de $y$ en $K_\bullet(X_s)$, podemos encontrar enteros $n_i$ tales que
\[
y_s=\sum_i\xi_i\cdot n_i.
\]
Por otra parte, se tiene $K=\varinjlim_f A_f$, donde $f$ recorre el conjunto de los elementos de $A$. Por tanto, $\Spec(K)=\varprojlim_f\Spec(A_f)$, y
\[
X_s=X\times_S\Spec(K)=\varprojlim_f X\times_S\Spec(A_f),
\]
por EGA IV 8.2.3 y 8.2.5. Como los esquemas considerados son noetherianos y los morfismos de transición afines y planos, se tiene por IV 3.2.3
\[
K_\bullet(X_s)=\varinjlim_f K_\bullet(X\times_S\Spec(A_f)).
\]
Como los homomorfismos de «~imagen inversa~» sobre los $K_\bullet$ son lineales respecto de los $K^\circ$ (IV (2.12.2)), existe un abierto no vacío $U=\Spec(A_f)$ de $S$ tal que en $K_\bullet(X_U)$ se tenga
\[
y_U=\sum_i\xi_i\cdot n_i.
\]
Sea $S'$ un subesquema cerrado de $S$ con soporte $S-U$. Entonces $\varphi_{S'}$ es sobreyectivo por la hipótesis de inducción, y por consiguiente el diagrama (1.2.1) muestra que
\[
y-\sum_i\xi_i\cdot n_i\in \operatorname{Im}(\varphi),
\]
por lo que, como $n_i=f^*(n_i)$, $\varphi$ es sobreyectivo.

\textbf{Corolario 1.3.} Sea $f:X\to S$ un morfismo plano de esquemas noetherianos. Se supone que, para todo punto $s$ de $S$, el homomorfismo
\[
K^\circ(X)\longrightarrow K_\bullet(X_s)
\]
es sobreyectivo. Entonces el homomorfismo deducido de $f^*$
\[
K_\bullet(S)\otimes_{K^\circ(S)}K^\circ(X)\longrightarrow K_\bullet(X)
\]
es sobreyectivo.

Tomamos como familia $(\xi_i)_{i\in I}$ la familia de todos los elementos de $K^\circ(X)$. Se verifican entonces las hipótesis de 1.2 y, por consiguiente, el homomorfismo
\[
\varphi:\bigl(K_\bullet(S)\bigr)^{(K^\circ(X))}\longrightarrow K_\bullet(X)
\]
es sobreyectivo. Por la definición de $\varphi$, esto equivale a decir que el homomorfismo
\[
K_\bullet(S)\otimes_{K^\circ(S)}K^\circ(X)\longrightarrow K_\bullet(X)
\]
es sobreyectivo.

\textbf{Corolario 1.4.} Sea $f:X\to S$ un morfismo plano de esquemas noetherianos. Se supone que, para todo punto $s$ de $S$, el homomorfismo canónico $\mathbb Z\to K_\bullet(X_s)$ es un isomorfismo. Entonces el homomorfismo
\[
f^*:K_\bullet(S)\longrightarrow K_\bullet(X)
\]
es sobreyectivo. Si, además, $f$ posee una sección $g$ de dimensión Tor finita, entonces $f^*$ es un isomorfismo.

Tomamos como familia $(\xi_i)_{i\in I}$ la familia reducida al único elemento $1\in K^\circ(X)$; se tiene entonces $\varphi=f^*$ y, como se cumplen las condiciones de 1.2, $f^*$ es sobreyectivo.

Si $f$ admite una sección $g$ de dimensión Tor finita, se puede definir un homomorfismo $g^*:K_\bullet(X)\to K_\bullet(S)$ (IV 2.12). Como se tiene $g^*\circ f^*=(f\circ g)^*=\operatorname{Id}$ por IV 2.12.1.1, $f^*$ es inyectivo y, por tanto, un isomorfismo.

\textbf{Observación 1.5.} Para que $g$ sea de dimensión Tor finita, es necesario y suficiente que las fibras de $X$ sean regulares en los puntos de $g(S)$, es decir, que $f$ sea regular en los puntos de $g(S)$ (o también liso en esos puntos si se supone que $f$ es de presentación finita).

Como estas propiedades son locales sobre $X$ y $S$, podemos suponerlos afines; sean $S=\Spec(A)$, $X=\Spec(B)$. Se tiene entonces:
\[
\operatorname{tor.dim}_B(A)=\operatorname{dim.proj}_B(A)=\sup_{s\in S}\operatorname{dim.proj}_{B\otimes_A k(s)}(k(s))
\]
por EGA IV 12.3.2, donde se puede sustituir la hipótesis de presentación finita sobre $B$ por hipótesis noetherianas. Ahora bien, $k(s)$ es isomorfo al cuerpo residual de $\mathcal O_{X_s,g(s)}$, y por EGA $0_{\mathrm{IV}}$ 17.2.5 se tiene:
\[
\operatorname{dim.proj}_{B\otimes_A k(s)}(k(s))=
\operatorname{dim.proj}_{\mathcal O_{X_s,g(s)}}(k(s)).
\]
Resulta entonces de EGA $0_{\mathrm{IV}}$ 17.2.11 que esta última es finita si y solo si $\mathcal O_{X_s,g(s)}$ tiene dimensión cohomológica finita, es decir, es regular.

\textbf{Proposición 1.6.} Sean $S$ un esquema noetheriano, $E$ un $\mathcal O_S$-módulo localmente libre de tipo finito, $X$ el fibrado vectorial asociado a $E$, y $f$ el morfismo estructural de $X$. Entonces \(f^*:K_\bullet(S)\longrightarrow K_\bullet(X)\) es un isomorfismo.

Como $f$ es liso y admite una sección, a saber, la sección cero, para demostrar 1.6 basta mostrar que, para todo $s\in S$, $K_\bullet(X_s)\simeq \mathbb Z$ (1.4). Quedamos así reducidos al caso en que el esquema base es un cuerpo.

En este caso, $E$ es suma directa de submódulos de rango $1$, lo que permite descomponer $f$ en una sucesión de fibraciones de rango $1$, y por tanto reducirse al caso en que $X$ es de rango $1$, es decir, $X=\Spec(k[T])$. La proposición resulta entonces del lema siguiente:

\textbf{Lema 1.7.} Sean $A$ un anillo principal, $X=\Spec(A)$. Entonces el homomorfismo canónico
\[
\mathbb Z\longrightarrow K_\bullet(X)
\]
es un isomorfismo.

Como $X$ es regular, el homomorfismo $K^\circ(X)\to K_\bullet(X)$ es un isomorfismo. Por tanto, $K_\bullet(X)$ está engendrado por las clases de los $\mathcal O_X$-módulos localmente libres de tipo finito, es decir, por las clases de los $A$-módulos proyectivos de tipo finito. Como $A$ es principal, todo módulo proyectivo de tipo finito es libre, de donde resulta la sobreyectividad. La inyectividad se debe a que $K^\circ(X)$ está $\mathbb Z$-aumentado.

\textbf{Proposición 1.8.} Sean $S$ un esquema noetheriano, $E$ un $\mathcal O_S$-módulo localmente libre de rango $r$, $X$ el fibrado vectorial asociado a $E$, $f$ el morfismo estructural de $X$, y $g$ la sección cero de $X$. Entonces el homomorfismo compuesto
\[
K_\bullet(S)\xrightarrow{g_*}K_\bullet(X)\xrightarrow{g^*}K_\bullet(S)
\]
(donde $g^*$ se define según IV 2.12.1 y 1.5) es la multiplicación por el elemento $\lambda_{-1}(E)$ de $K^\circ(S)$, denotando también por $E$ la clase de $E$ en $K^\circ(S)$, y poniendo
\[
\lambda_{-1}(E)=\sum_{i=0}^r(-1)^i\lambda^i(E).
\]

Este enunciado es el análogo, para $K_\bullet$, de VII 2.7. Sea $F$ un $\mathcal O_S$-módulo coherente, de clase $F$ en $K_\bullet(S)$. Entonces $g_*(F)$ es la clase en $K_\bullet(X)$ del complejo $\mathbb Rg_*(F)=g_*(F)$, y por tanto de $F$ considerado como $\mathcal O_X$-módulo por medio de $g$. Para calcular $g^*(g_*(\sisi{f}{F}))$, hay que tomar la clase de \(\mathbb Lg^*(F)=F\otimes^{\mathbb L}_{\mathcal O_X}\mathcal O_S\).
Para ello, consideramos el complejo
\[
0\longrightarrow \Lambda^r(E)\otimes_{\mathcal O_S}S(E)\longrightarrow
\Lambda^{r-1}(E)\otimes_{\mathcal O_S}S(E)\longrightarrow\cdots\longrightarrow
E\otimes_{\mathcal O_S}\mathbb S(E)\longrightarrow \mathbb S(E)\longrightarrow 0,
\]
donde $S(E)$ es el álgebra simétrica de $E$. Es una resolución de $\mathcal O_S$ sobre $S(E)$ (VI 1.11). El complejo asociado sobre $X$, denotado $L_\bullet$, tiene como término $k$-ésimo $\Lambda^k(f^*(E))$ y es, por tanto, una resolución de $\mathcal O_S$ por $\mathcal O_X$-módulos localmente libres de tipo finito; se tiene así \(L_\bullet\otimes_{\mathcal O_X}F=\mathcal O_S\otimes^{\mathbb L}_{\mathcal O_X}F\). Por consiguiente:
\[
\begin{aligned}
g^*(g_*(F))
&=\cl(L_\bullet\otimes_{\mathcal O_X}F)
=\sum_{i=0}^r(-1)^i\cl(L_i\otimes_{\mathcal O_X}F)\\
&=\sum_{i=0}^r(-1)^i\cl(\Lambda^i(E)\otimes_{\mathcal O_S}F)
=F\cdot\lambda_{-1}(E).
\end{aligned}
\]
La proposición se sigue por aditividad.

